% ==============================================================================
\section{MILP Formulation}
\label{sec:milp}
% ==============================================================================

We now present a mixed-integer linear programming formulation of the AHPP.
The model is implemented in Pyomo~\citep{Bynum2021} and solved with
Gurobi~\citep{Gurobi2023}.  Throughout this section parameters are written
with upper-case Latin letters and decision variables with lower-case Latin
letters, with Greek letters reserved for fixed scalar constants such as the
minimum separation $\varepsilon$ and the big-$M$ relaxation $M$.

\subsection*{Sets}

\begin{tabular}{@{}l@{\quad}p{0.78\textwidth}@{}}
  $\calP$
    & maintenance positions \\
  $\calR$
    & aircraft \\
  $\calJ$
    & jobs; $\calJ_r \subseteq \calJ$ denotes the jobs of aircraft $r$ \\
  $\calA \subseteq \calP \times \calP$
    & blocking arcs $(p, p')$: position $p$ blocks position $p'$ \\
  $\calE \subseteq \calJ \times \calJ$
    & job precedences $(j, j')$: job $j$ must finish before $j'$ starts \\
\end{tabular}

\subsection*{Parameters}

\begin{tabular}{@{}l@{\quad}p{0.78\textwidth}@{}}
  $D_j$
    & processing time of job $j$ \\
  $E_r$
    & earliest start time of aircraft $r$ \\
  $L_r$
    & target finish time (due date) of aircraft $r$ \\
  $R_{jr}$
    & $1$ if job $j$ belongs to aircraft $r$, $0$ otherwise \\
  $A_{jr}$
    & $1$ if job $j$ is the first job of aircraft $r$, $0$ otherwise \\
  $\Omega_{jr}$
    & $1$ if job $j$ is the last job of aircraft $r$, $0$ otherwise \\
  $\varepsilon$
    & tow-in / tow-out time, used both as separation between consecutive
      aircraft at the same position and as the safety margin
      $f_r - \varepsilon$ in the blocking conditions; per-instance value
      read from the input \\
  $H^{UB}$
    & technical upper bound on every aircraft start and finish time,
      $H^{UB} = \max_{r \in \calR} E_r + \sum_{r \in \calR} D_r +
      (|\calR| - 1)\varepsilon$ \\
  $M$
    & big-$M$ constant; we take $M = H^{UB} + \varepsilon$ \\
  $W^M, W^D, W^S$
    & weights for makespan, total delay, and movements \\
\end{tabular}

\smallskip
\noindent
The binary parameters $R_{jr}, A_{jr}, \Omega_{jr}$ are assumed
consistent: every job belongs to exactly one aircraft
($\sum_r R_{jr} = 1$ for all $j$), every aircraft has exactly one first
and one last job ($\sum_j A_{jr} = \sum_j \Omega_{jr} = 1$ for all $r$),
and first/last jobs of $r$ must belong to $r$ ($A_{jr}, \Omega_{jr} \le
R_{jr}$).

\subsection*{Variables}

\paragraph{Scheduling and assignment.}
Continuous variables $t^S_j, t^F_j \in \NR$ are the start and finish
times of job $j$, and $s_r, f_r \in \NR$ are the start and finish times
of aircraft $r$ at its assigned position.  The binary variables are
$x_{rp}$ (1 if aircraft $r$ is assigned to position $p$),
$y_{jp}$ (1 if job $j$ is assigned to $p$), and $w_{jj'p}$ (1 if job $j$
precedes job $j'$ at position $p$).

\paragraph{Blocking.}
For every pair of distinct aircraft $r \ne r'$ and every blocking arc
$(p, p') \in \calA$, we introduce six binary variables.  On the entry
side, $\alpha^{\mathrm{in}}_{rr'pp'} = 1$ iff $s_{r'} > s_r$,
$\beta^{\mathrm{in}}_{rr'pp'} = 1$ iff $s_{r'} \le f_r - \varepsilon$, and
$u^{\mathrm{in}}_{rr'pp'} = 1$ iff the entry of $r'$ at $p'$ forces
displacing $r$ at $p$.  The exit-side variables
$\alpha^{\mathrm{out}}_{rr'pp'}, \beta^{\mathrm{out}}_{rr'pp'},
u^{\mathrm{out}}_{rr'pp'}$ are defined symmetrically by replacing $s_{r'}$
with $f_{r'}$ in the conditions above.

\paragraph{Objective.}
The makespan $m \in \NR$, the per-aircraft delay $v^D_r \in \NR$, and the
total number of movements $n \in \NZ$ collect the three components of the
objective.

\subsection*{Constraints}

To keep the constraints compact, we abbreviate the universal quantifiers
$\forall j \ne j' \in \calJ,\, p \in \calP$ by $(\ast)$ and
$\forall r \ne r' \in \calR,\, (p, p') \in \calA$ by $(\dagger)$.
The full model is:
\begin{align}
  \sum_{p \in \calP} x_{rp} &= 1
    && \forall r \in \calR \label{eq:assign}\\
  y_{jp} &= x_{rp}
    && \forall j \in \calJ,\, r \in \calR,\, p \in \calP : R_{jr} = 1
    \label{eq:linkage}\\
  t^F_j &= t^S_j + D_j
    && \forall j \in \calJ \label{eq:finish}\\
  s_r &= \sum_{j \in \calJ} A_{jr}\, t^S_j
    && \forall r \in \calR \label{eq:srdef}\\
  f_r &= \sum_{j \in \calJ} \Omega_{jr}\, t^F_j
    && \forall r \in \calR \label{eq:frdef}\\
  t^S_j &\ge E_r
    && \forall j \in \calJ,\, r \in \calR : A_{jr} = 1
    \label{eq:earlystart}\\
  t^S_{j'} &\ge t^F_j
    && \forall (j, j') \in \calE \label{eq:prec}\\
  t^S_{j'} &\ge t^F_j
              - M\bigl(1 - w_{jj'p} + (1 - y_{jp}) + (1 - y_{j'p})\bigr)
    && (\ast) \label{eq:nonov1}\\
  t^S_{j}  &\ge t^F_{j'}
              - M\bigl(w_{jj'p} + (1 - y_{jp}) + (1 - y_{j'p})\bigr)
    && (\ast) \label{eq:nonov2}\\
  w_{jj'p} + w_{j'jp} &\ge y_{jp} + y_{j'p} - 1
    && (\ast) \label{eq:nonov3}\\
  w_{jj'p} + w_{j'jp} &\le 1
    && (\ast) \label{eq:nonov4}\\
  \sum_{j, j' \in \calJ}\!
    \bigl(\Omega_{jr} A_{j'r'} w_{jj'p}
        &+ \Omega_{j'r'} A_{jr} w_{j'jp}\bigr) \notag\\
    &\ge x_{rp} + x_{r'p} - 1
    && \forall r < r' \in \calR,\, p \in \calP \label{eq:nonov-ac}\\
  s_{r'} &\ge s_r - M\bigl(1 - \alpha^{\mathrm{in}}_{rr'pp'}\bigr)
    && (\dagger) \label{eq:alphaInLB}\\
  s_{r'} &\le s_r + M\,\alpha^{\mathrm{in}}_{rr'pp'}
    && (\dagger) \label{eq:alphaInUB}\\
  s_{r'} &\le f_r - \varepsilon
              + M\bigl(1 - \beta^{\mathrm{in}}_{rr'pp'}\bigr)
    && (\dagger) \label{eq:betaInUB}\\
  s_{r'} &\ge f_r - M\,\beta^{\mathrm{in}}_{rr'pp'} \notag\\
         &\quad - M(1 - x_{rp}) - M(1 - x_{r'p'})
    && (\dagger) \label{eq:betaInLB}\\
  u^{\mathrm{in}}_{rr'pp'} &\le \min\{x_{rp},\, x_{r'p'},\,
                                   \alpha^{\mathrm{in}}_{rr'pp'},\,
                                   \beta^{\mathrm{in}}_{rr'pp'}\}
    && (\dagger) \label{eq:uInUB}\\
  u^{\mathrm{in}}_{rr'pp'} &\ge x_{rp} + x_{r'p'} \notag\\
                           &\quad + \alpha^{\mathrm{in}}_{rr'pp'}
                           + \beta^{\mathrm{in}}_{rr'pp'} - 3
    && (\dagger) \label{eq:uInLB}\\
  f_{r'} &\ge s_r - M\bigl(1 - \alpha^{\mathrm{out}}_{rr'pp'}\bigr)
    && (\dagger) \label{eq:alphaOutLB}\\
  f_{r'} &\le s_r + M\,\alpha^{\mathrm{out}}_{rr'pp'}
    && (\dagger) \label{eq:alphaOutUB}\\
  f_{r'} &\le f_r - \varepsilon
              + M\bigl(1 - \beta^{\mathrm{out}}_{rr'pp'}\bigr)
    && (\dagger) \label{eq:betaOutUB}\\
  f_{r'} &\ge f_r - M\,\beta^{\mathrm{out}}_{rr'pp'} \notag\\
         &\quad - M(1 - x_{rp}) - M(1 - x_{r'p'})
    && (\dagger) \label{eq:betaOutLB}\\
  u^{\mathrm{out}}_{rr'pp'} &\le \min\{x_{rp},\, x_{r'p'},\,
                                    \alpha^{\mathrm{out}}_{rr'pp'},\,
                                    \beta^{\mathrm{out}}_{rr'pp'}\}
    && (\dagger) \label{eq:uOutUB}\\
  u^{\mathrm{out}}_{rr'pp'} &\ge x_{rp} + x_{r'p'} \notag\\
                            &\quad + \alpha^{\mathrm{out}}_{rr'pp'}
                            + \beta^{\mathrm{out}}_{rr'pp'} - 3
    && (\dagger) \label{eq:uOut}\\
  n &= 2 \!\!\sum_{(p,p') \in \calA}
         \sum_{\substack{r, r' \in \calR\\ r \ne r'}}
         \bigl(u^{\mathrm{in}}_{rr'pp'} + u^{\mathrm{out}}_{rr'pp'}\bigr)
    && \label{eq:movements}\\
  v^D_r &\ge \sum_{j \in \calJ} \Omega_{jr}\, t^F_j - L_r
    && \forall r \in \calR \label{eq:delay}\\
  m &\ge t^F_j
    && \forall j \in \calJ \label{eq:makespan}
\end{align}

Constraint~\eqref{eq:assign} forces each aircraft to occupy exactly one
position, and \eqref{eq:linkage} propagates this assignment from aircraft
to their jobs through the link parameter $R_{jr}$; the duration of each
job is encoded in \eqref{eq:finish}, while
\eqref{eq:srdef}--\eqref{eq:frdef} collapse the start and finish times of
each aircraft onto its first and last job using the parameters
$A_{jr}$ and $\Omega_{jr}$, and \eqref{eq:earlystart} pushes the
aircraft's earliest-start window onto that first job.  Internal job
precedences are imposed by \eqref{eq:prec}.  Non-overlap at a position
is enforced in two layers: at the job level,
\eqref{eq:nonov1}--\eqref{eq:nonov2} use the ordering binary $w_{jj'p}$
to choose the disjunction direction, while
\eqref{eq:nonov3}--\eqref{eq:nonov4} guarantee that whenever two jobs
share a position exactly one direction is active and that both cannot
hold simultaneously; at the aircraft level, \eqref{eq:nonov-ac} closes
the gap by requiring that whenever two aircraft share a position the
last job of one is ordered before the first job of the other.  The
entry-blocking block defines two binary indicators and combines them:
\eqref{eq:alphaInLB}--\eqref{eq:alphaInUB} make
$\alpha^{\mathrm{in}}_{rr'pp'}$ equal to $1$ iff $s_{r'} > s_r$, while
\eqref{eq:betaInUB}--\eqref{eq:betaInLB} make
$\beta^{\mathrm{in}}_{rr'pp'}$ equal to $1$ iff $s_{r'}$ falls in the
servicing window of $r$, the position guards
$-M(1 - x_{rp}) - M(1 - x_{r'p'})$ in \eqref{eq:betaInLB} switching the
constraint off whenever the blocking placement does not actually hold;
\eqref{eq:uInUB}--\eqref{eq:uInLB} then set
$u^{\mathrm{in}}_{rr'pp'}$ to the conjunction
$x_{rp} \wedge x_{r'p'} \wedge \alpha^{\mathrm{in}}_{rr'pp'} \wedge
\beta^{\mathrm{in}}_{rr'pp'}$, so an entry movement is counted iff both
aircraft are at the relevant blocking positions and $r'$ enters during
the servicing window of $r$.  The exit-blocking block
\eqref{eq:alphaOutLB}--\eqref{eq:uOut} mirrors the entry block with
$f_{r'}$ in place of $s_{r'}$.  Finally, \eqref{eq:movements} aggregates
all entry and exit indicators into the movement count $n$, with the
factor $2$ accounting for displacement plus return;
\eqref{eq:delay} lower-bounds the positive delay of each aircraft by
$f_r - L_r$; and \eqref{eq:makespan} lower-bounds the makespan by every
job finish time.

\subsection*{Objective}
\begin{equation}
  \min \quad W^M\, m
         \;+\; W^D \sum_{r \in \calR} v^D_r
         \;+\; W^S\, n.
  \label{eq:milpobj}
\end{equation}

